% Source: https://artofproblemsolving.com/community/c3h1874653_supplementary_angles
% Stable identity: aops-topic-1874653-q01
Let the smaller angle be $x$. Supplementary angles add up to $180^\circ$ (together they form a straight angle):

\begin{asy}
usepackage("amsmath");
size(6cm);
defaultpen(fontsize(10pt));
pair O=(0,0);
draw((-3,0)--(3,0),linewidth(1));draw(O--2.8*dir(18),linewidth(1));dot(O);
draw(arc(O,1.3,0,18),blue);label("$x=18^\circ$",1.3*dir(9),E,blue);
draw(arc(O,0.7,18,180),red);label("$9x=162^\circ$",0.7*dir(100),N,red);
label("$O$",O,S);
label("\mbox{}",(min(currentpicture,user=true)+max(currentpicture,user=true))/2,black); // keeps the text after the figure black when LaTeX typesets the labels inline
\end{asy}

\[
x+9x=180^\circ\implies x=18^\circ,
\]
so the larger angle is $9\cdot18^\circ=\boxed{162^\circ}$.
